\section{Five permanent seven-plane obstructions}

\begin{finding}{P-13}{five permanent coordinate orbits have border rank at least fourteen}
In the standard rectangle numbering, the permanent residual tensors
for each of the following subsets have border rank at least fourteen
over every algebraically closed field of characteristic zero:
\[
\begin{aligned}
 C_1&=\{0,1,2,3,4,6,12\},&
 C_2&=\{0,1,2,3,6,7,8\},\\
 C_3&=\{0,1,2,6,7,8,18\},&
 C_4&=\{0,1,2,6,7,18,20\},\\
 C_5&=\{0,1,3,6,7,9,14\}.&&
\end{aligned}
\]
This is an assertion about five specified orbits; it does not alone
evaluate the uniform seven-plane minimum.
\end{finding}
\begin{proof}
Each tensor has format $7\times12\times12$, and its last-factor
image $U\subset A\otimes B$ has dimension twelve. If its border
rank were at most thirteen, the necessary conditions preceding
R-04 would supply a thirteen-plane $L\supset U$ with
\[
 \rank\Psi_A(L)\leq78,\qquad\rank\Psi_B(L)\leq143.
\]
The row/column torus preserves the tensor and these determinantal
incidences. The one-parameter fixed-limit argument of P-09 therefore
allows $L/U$ to be one homogeneous line.

We use exactly the complete saturation and linear-column criterion
proved in P-09, with permanent coefficients $+1$ in all four
ordered entries of each slice. The finite certificate gives:
\begin{center}
\begin{tabular}{c|r|r|c|r|r}
Case & Quotient weights & Saturation exclusions & Remaining map
 & Baseline rank & Required increment\\\hline
$C_1$ &60&58&$\Psi_A$&74&5\\
$C_2$ &52&49&$\Psi_A$&74&5\\
$C_3$ &50&45&$\Psi_B$&136&8\\
$C_4$ &50&46&$\Psi_A$&74&5\\
$C_5$ &56&56&none&78&--
\end{tabular}
\end{center}
For every saturation exclusion, the certificate exhibits a nonzero
minor of order $\dim(L_w\otimes F)-12$ in an exterior map on
$L_w=U+V_w$. Its kernel has dimension at most twelve. Inclusion
$L\subset L_w$ then contradicts the necessary kernel dimension
at least thirteen, as in P-09.

For $C_1,C_2,C_4$ the remaining quotient blocks number two,
three and four, respectively. An integer annihilator of each
baseline exterior image turns vanishing of a new projected
column into explicit linear conditions on a nonzero quotient
vector. Every stack for three of the seven $A$ columns has
full column rank. Thus at most two columns vanish, and the
rank increment is at least five. Its distinct nonzero column
weights ensure independence. Hence $\rank\Psi_A(L)\geq79>78$.
For $C_3$ there are five remaining blocks. Every stack for
five of the twelve $B$ columns has full column rank, so at
least eight projected columns are nonzero and independent.
Thus $\rank\Psi_B(L)\geq144>143$.
Case $C_5$ has no remaining block. All possibilities have
therefore been excluded.

The data in \path{per7_m12_certificate.json} specify the
baseline and saturation minors and all integer annihilators.
The independent \path{verify_seven_plane_certificate.py}
reconstructs the ordered permanent coefficients and exterior
maps directly; it imports no generator or rank/nullspace oracle.
It verifies all $268$ quotient-weight choices, $254$ saturation
minors, and $4275$ linear zero-pattern tests. All annihilator
identities are checked over the integers, and all claimed
nonzero minors are checked modulo the verified prime $1000003$.
The integer products obey the same explicit bounds as in P-09.
A nonzero residue implies a nonzero integer minor, proving
the stated characteristic-zero assertions.
\end{proof}

The geometric dictionary is P-09's proper Grassmannian incidence,
quotient weight bundles and empty projective linear loci, with the
permanent sign convention. It does not require a permanent Borel
compression, and none is asserted.
